\chapter{114}

\section{Mailand (IT), im November}

\subsection{Louis}

Wie genau Joh alles wissen will.

Ich mag seinen ruhigen Fahrstil. Der Verkehr fliesst. Wir kommen
ungestört voran. In wenigen Minuten erreichen wir die Gotthard
Raststätte.

Es ist, als sei ich noch einmal mitten drin, während ich von all den
Begegnungen erzähle. Meine Reise nach Zürich und spätestens das
Zusammentreffen mit Jess lässt ihn aufhorchen. Und wie ich von Mascha
berichte, wird er unruhig. Zur Beschaffung meines falschen Passes stellt
er mir einige Fragen. Ich spüre, er hält sich mit Urteilen zurück. Ich
bin froh darüber.

«Unglaublich, Louis, dass unser Inserat offen auf dem Tisch dieser
Familie lag, wirklich nicht zu fassen.»

Für einen kurzen Augenblick dreht er den Kopf zu mir. Ich verschränke
die Finger hinter meinem Hals und atme durch.

«Ja. Mir kam euer Aufruf vor, als wäre er an mich gerichtet. Ich war so
verzweifelt. Bedenken hatte ich schon. Ich hatte natürlich keine Ahnung,
was mich bei euch erwarten würde,» ich nehme meine Arme runter und für
einen Moment wird mir eng, «und hatte keinen Plan, was ich tun sollte,
wenn es nicht klappen würde.»

Joh nimmt die Ausfahrt auf die Raststätte. Die letzte auf Schweizer
Boden.

Meine Beine kribbeln.

«Phu, gut, sind wir auf der andern Seite.»

Joh löst den Sicherheitsgurt und scheint keine Eile zu haben,
auszusteigen.

Er dreht sich zu mir.

«Hm, da drüben also stiegst du zu Leonard um?»

«Ja, und verabschiedete mich von Giovanni.»

Ich hänge meinen Gedanken nach. Sehe Giovanni vor mir, sein fröhliches
Lachen und habe das Gefühl, dass es schon ewig her ist, seit er mich
gerettet hat. Und noch heute nichts davon weiss. Gerade bin ich ihm sehr
nah.

Johs Handy unterbricht uns.

«Ja, B.B.?»

Joh hört konzentriert zu. Dann antwortet er. Ben scheint sich nach einem
bestimmten Ablauf zu erkundigen.

Mein Blick bleibt auf Johs braungebrannter, breiter Hand liegen. Ich
widerstehe dem Impuls, sie zu berühren. Vielleicht sollte ich aufhören,
mich dauernd zu bedanken. Ich glaube, er weiss wirklich, wie sehr ich
seine Unterstützung schätze.

\sectionbreak

Sie hatten Johs Plan alle vor mir gekannt.

Meine Worte hatten bestimmt vorwurfsvoll geklungen. \emph{«Du hast mir
nichts gesagt.»}

Joh musste bemerkt haben, wie eingeschnappt ich war. «Aber natürlich
nicht, Louis,» hatte er mir erzählt, «bevor ich dir das Angebot
unterbreiten konnte, musste ich sicher stellen, dass all meine Aufgaben
auf der «Sunnewaid» auf die andern verteilt werden konnten. Ja, und vor
allem, dass sie meine Idee überhaupt unterstützten. Erst nachdem ich Ben
im Boot hatte und er obendrauf als Stellvertreter gegenüber David
gutgeheissen wurde, war der Weg offen.»

Seine Idee der Rückwärtsverfolgung meiner Reise, wie er es nannte, hatte
mich erstmal erschreckt. All diesen Orten der Verzweiflung noch einmal
zu begegnen? Konnte das gutgehen? Und, was hatte ich davon? Joh hatte
mich mit keinem einzigen Wort versucht, zu überzeugen. Auch auf meine
Frage, was er sich davon verspreche, hatte er nur mit den Schultern
gezuckt.

Noch sehe ich mich ratlos oben auf der Bank sitzen. Meine Entscheidung
war erst am nächsten Morgen gekommen. Dafür glasklar.

\sectionbreak

Wir halten unsere Pause kurz. Joh ist meine Unruhe aufgefallen. Ich kann
meine Beine kaum mehr still halten.

Er stellt unser Geschirr zusammen. Dann schaut er mich liebevoll an.

«Komm, Louis, weiter geht’s.»

Jetzt bin ich dran. Wie wechseln uns beim Fahren ab. Ich steige mit
Vorfreude ein. Endlich komme ich wieder einmal dazu, länger am Steuer zu
sitzen. Unser Mietwagen hat Power. Ich fühle mich ziemlich smart.

Joh scheint meinen Fahrkünsten zu vertrauen. Er hat seinen Sitz schräg
gestellt und die Mütze über die Augen gezogen.

Der kühle Novembertag unterscheidet sich ungemein vom Tag der Hinfahrt,
diesem heissen, hellen und in jeder Hinsicht schrecklichen. Die
vorbeiflitzende Landschaft präsentiert sich gedeckt und bereitet sich
unabwendbar auf den Winter vor. Was Giovanni wohl gerade tut? Seine
Perry Como- Geschichte fällt mir ein. Und meine Erinnerung, wie wir uns
zusammen \emph{«Magic Moments»} angehört haben.

«Joh?»

Ein knappes Grunzen erreicht mich vom Beifahrersitz.

«Stört es dich, wenn ich Musik anstelle?»

«Nö.»

«Okay. Giovanni hat während der Fahrt Stücke von Perry Como laufen
lassen. \emph{«Magic Moments»} war eines davon.»

Ich brauche ein bisschen, bis ich zwischen Gas geben, Überholen und
Bremsen den Song gefunden habe. Mit den ersten Tönen hebt Joh seine Arme
und macht dirigierende Bewegungen. Unter der Mütze sehe ich knapp seinen
lachenden Mund.

Mir ist, als habe ich andere Ohren auf als in Giovannis Laster. Der Song
passt heute um einiges besser. Ich fühle mich leicht.

\qrblock{Perry Como \emph{«Magic Moments»}}{media/QR-Codes/012_Magic_Moments_Perry_Como_Spotify.png}

\subsection{Louis}

Wir sind geschafft. Alle beide.

Ich fühle mich, als komme ich nach Jahren zurück in meine Wohnung. Ich
öffne die Haustür und lasse Joh vor mir eintreten. Er stellt seine
Tasche im schmalen Flur ab und schaut sich interessiert um.

«Hübsch hast du es hier, Louis.»

Trotz Mamans gelegentlichem Lüften riecht es muffig. Immerhin ist es
angenehm warm. Wie angeboten, hat sie die Heizung angedreht. Als Erstes
ziehe ich sämtliche Rollläden hoch und öffne die Fenster.

«Ja, klein zwar,» rufe ich in seine Richtung, «aber ganz okay.»

Mit Blick ins Schlafzimmer sehe ich meinen Gitarrenkoffer ans Bett
gelehnt stehen. Ich nähere mich ihm, hebe ihn hoch und öffne den Deckel
mit wackligen Händen. Verloren und alleingelassen steht sie aufrecht
drin, meine \emph{«Taylor»}. Mir kommt vor, als habe sie einen
vorwurfsvollen Ausdruck. Valentina musste Maman die Gitarre gebracht
haben. Es fällt mir gar nicht ein, dass ich mich beherrschen könnte. Ich
sehe den Koffer auf dem \emph{Roccia} vor dem grossen Sitzstein stehen.
Als letztes Bild von ihm vor meiner Flucht. Meine Augen füllen sich
stossweise mit Tränen.

Ich höre die Klospülung und Joh die Türe der Toilette schliessen. Beim
Näherkommen wandert sein Blick von mir zur Gitarre und bleibt zurück auf
mir liegen. Er bleibt neben mir stehen, bückt sich und legt mir seine
Hand sanft an den Hals. Dann zieht er meinen Kopf auf seine Schulter.

\sectionbreak

Wir haben den Ablauf der kommenden Tage mehrmals durchgespielt. Für
heute ist alles klar. Selbst das Lebensmittelgeschäft um die Ecke konnte
ich Joh beim Vorbeifahren zeigen. Er macht einen erschöpften Eindruck.

«Ich hoffe, du kannst dich etwas ausruhen, Joh, bis wir da sind.»

«Ja, ja, passt schon. Weder der Einkauf noch die Zubereitung des
Nachtessens wird lange dauern. Mach dir keine Gedanken. Um sieben werde
ich euch mit einem wahren Dinner überraschen. Wann willst du bei deiner
Mutter sein?»

«Maman hat sich den Nachmittag frei gehalten. Ich dusche kurz und geh
dann auch gleich los. Es zieht mich zu ihr. Und, ich freue mich wie ein
Kind auf unseren gemeinsamen Abend.»

Joh gähnt und hält sich die Hand vor den Mund.

«Ich mich auch, und ganz speziell darauf, deine Mutter kennenzulernen.»
